\chapter{User Study} \label{sec:UserStudy}

As a part of this research project, outreach was conducted to individuals in the
game playing\/development space. A survey was sent out to acquaintances of the
author, as well as members of Cal Poly Game Development Club, collecting information
on prior experience of the participants as well as information about how they used
the system. All survey materials were approved by the Cal Poly San Luis Obispo IRB
(IRB Protcol \#2026-080). There were 4 total responses to the final survey, and
below are some primary takeaways from the responses made by those who completed it.

\section{Engine Capabilities} \label{sec:EngineCapabilities}

Users seemed to be satisfied with the capabilities of the language and engine. Two
expressed this directly, stating that they feel like they could do much more with
it if given the time to explore. One mentioned that the modularity of the practice
system gave them many options for setting up their scenarios and dropping characters
into different situations.

As many critiques as there were, there appeared to be minimal concerns about the
core concepts explored with Praxsmith. All participants were satisfied with their
resulting product after experimentation. The challenge then becomes adjusting how
they interact with those inner systems so that they can convey their ideas more
easily.

\section{Documentation And Examples} \label{sec:DocumentationAndExamples}

As the developer of any project, it is hard to evaluate the approachability of a
system from the standpoint of someone who has never touched it. Great care was taken
in building a documentation page for the language, including tips on how to use it
and three short examples showcasing different features and potentials.

All participants had positive reviews about the documentation and examples, with
2\/4 participants rating the usefulness of the documentation at full score on the
likert scale and 3\/4 participants giving a full score to the provided examples.
One user pointed out in one of their responses that a question they had showed up
in one of the examples, allowing them to quickly move forward with their project.

As a contrary point, one user expressed confusion about the purpose of the
$delete\ self$ outcome found in many of the example practices, which is a clarity
issue with the documentation. This confusion proves that although the existence of the
documentation was much appreciated, there are still holes in it that should be
identified and patched by more user studies.

\section{More Types} \label{sec:MoreTypes}

As one of the users was experimenting with an inventory system within their game,
they expressed the desire for array type support within the language. This would allow
for more complex interaction, like characters holding several items at once or
remembering a series of actors for later actions.

It is worth noting that arrays are very distinct datatypes and would alter the shape
of the language significantly. It was for this reason that they were not included, and
situations such as actors holding many things can instead be represented by an
exclusive $is\_held\_by$ relation and by adding dummy objects into the world.

However, this complaint highlights a shortcoming of the language either way. Either
the existing solution needs to be made more obvious and easier to work with, or
arrays should be added into the system as a way to support this functionality.
Inventory systems are incredibly common in games, so more research into this area
would be appropriate.

\section{If-Statements In Effects} \label{sec:IfStatementsInEffects}

One user noted the lack of a proper system around comparisons within the language.
Praxsmith provides if-statements for boolean values in the form of $and$ short
circuiting and can be extended to if\/else by adding $or$ and $not$ statements, but
there are two main problems with this:

\begin{itemize}
    \item It is clunky. Especially once you get into if\/else territory, the long
    expressions quickly become unreadable.
    \item It only works when returning boolean values, which is okay for conditions
    but outcomes\/effects suffer as a result. Changing dialog slightly based on some
    state variable is one example of where this might be useful, as that currently
    requires one actions within a practice for each dialog choice, which gets
    repetitive quickly. One survey participant pointed out this repetition in particular
    as a point of inconvenience.
\end{itemize}

This is likely the first target for language design change, as it would improve quality
of life greatly and allow for much more dynamic and expressive practice outcomes.

\section{Unique Syntax} \label{sec:UniqueSyntax}

Users of the language were universally surprised by the syntax at first, but said they
were able to adjust quickly. One user pointed out practices (e.g. $practice.greet.alice.bob$)
as a core example, noting that it became intuitive after regular usage.

This syntax is carried over from the sentences outlined in the original Praxis paper. (TODO: cite)
Since this project no longer strictly follows the exclusion logic laid out in
Praxis, sentences could be replaced with a more appropriate representation. On the other
hand, they work nicely as a singular unit and may be the cleanest option, especially
considering users seemed to adjust quickly. More research should be put into this
possibility.

\section{User Interface} \label{sec:UserInterface}

Only two users mentioned minor complaints with the user interface of the project.
One was to do with the visual highlighting being unclear, and another expressed that
they would like to be able to see the scripts as they are playing the game for
reference. Given the minimal nature of these complaints, it is reasonable to assume
that the development environment did not significantly impact user experience with the language.

\section{Survey Bias} \label{sec:SurveyBias}

This survey came with a significant amount of selection bias, as all respondents
to the preliminary survey had high experience with programming languages,
likely due to the areas this project was advertised. An important part of
this project was how easy the system is to use for writers, which ended up not
being determinable from the selection of survey participants. More analysis of
this point can be found in \Cref{sec:WriterFriendliness}, but it is worth noting
that more trials should be run specifically with writers who are not familiar
with programming for a more holistic interpretation.
